% ============================================================

\safechapter{Revisión crítica de algunos detalles de los apuntes}

\safesection{Conversión de MPa}
En material docente puede aparecer la aproximación:
\[
1\ \mathrm{MPa}\approx10\ \mathrm{kp/cm^2}.
\]
Para cálculos precisos:
\[
\boxed{1\ \mathrm{MPa}=10.19716\ \mathrm{kp/cm^2}}.
\]
La expresión con 10 es simplemente un redondeo ingenieril.

\safesection{Equilibrio y tercera ley de Newton}
Acción y reacción explica que las acciones internas aparecen por parejas iguales y opuestas.

Sin embargo, la condición de equilibrio de un cuerpo se expresa formalmente mediante:
\[
\sum\vect F=\vect0,
\qquad
\sum\vect M=\vect0.
\]
Conviene no identificar ambas ideas como si fueran exactamente la misma ley.

\safesection{Unicidad}
La afirmación “la solución es única” debe entenderse bajo las hipótesis de un problema elástico correctamente planteado.

Si una estructura puede desplazarse libremente como sólido rígido, la solución en desplazamientos no queda fijada.

% ============================================================
